\section{Conclusion}

The experiment isolates the effect of stochastic gold dynamics while holding operations and corporate assumptions fixed. Under the refreshed 2 September inputs, all four conditional medians lie between USD~35.14 and USD~35.87, while the completed USD~44.13 market close remains above each median. This is a directional comparison, not a claim of market mispricing.

Model quality is multidimensional. Heston provides the strongest current-surface fit, whereas Full Bates--Hawkes produces the lowest date-equal rolling OOS loss. Its 0.5569 branching ratio is stationary but no longer economically negligible. None of the four models beats the interpolated persistence surface, although Full Bates--Hawkes comes closest. Team 4's updated operating layer uses Q1--Q2 2026 actuals and a separate Q3-onward forecast, but its uncertainty is frozen before model comparison.

The paper should therefore be read as a controlled attribution exercise. The explicit growth, WACC and terminal-value assumptions make the DCF auditable, but the terminal block still contributes about 78\% of positive enterprise value. Any fair-value interpretation would additionally require a reconciled commodity basis, physical-measure mapping, operating uncertainty and equity bridge.

\begin{findingbox}[Decision-relevant takeaway]
Under a common operating and DCF contract, changing the stochastic gold engine can affect both the central valuation signal and tail risk. The current architecture is informative for model comparison, but it is not yet a filing-reconciled fair-value model.
\end{findingbox}
